%ez az én fájlom Szabi chill


\subsubsection{Általunk megvalósított modell}

A memrisztor szimulációjához használt .sub (subcircuit) fájl definiálja a saját komponensként implementált memrisztort. A saját modell "plus" és "minus" kapcsai között vannak definiálva a modellparaméterek, majd az ezekkel leírt differenciálegyenletek [melyikek?!]
A modellparaméterek:
\begin{itemize}
	\setlength\itemsep{0.35em}
	\item $R_\text{on}$ az erős adalékolású régió ellenállása
	\item $R_\text{off}$ a közel szigetelő (oxid) régió ellenállása
	\item $R_\text{ini}$ egy kezdőérték a szimuláció elindításához, a memrisztancia t=0 pillanatbeli értéke
	\item $\mu_v$ a töltéshordozók $\mu$ mozgékonysága (10f értéke a "femto" azaz 10-15 SI prefixum tízszerese)
	\item $D$ az eszköz teljes hossza
	\item $k$ és a pedig konstansok az előző paraméterek átláthatóbb használata érdekében bevezetve
\end{itemize}
\begin{figure}[H]
	\centering
	\includegraphics[scale=0.3]{memristor_sine_schematic}
	\caption{Általunk megvalósított ideális memrisztor modell és kapcsolás (schematic)\ref{appendix:kod}}
\end{figure}
\textit{MATLAB és LTSpice modellek egyeztetése}: MATLAB programkóddal is implementálhatjuk ugyanazon modellt. A függelékben csatolt kódban ugyanazon paraméterek lettek beállítva. 
A memrisztor kezdeti memrisztanciája [] alapján, $x =0.2$ estén $12.82 \kOhm$. 
A kapcsokra pedig egy feszültségforrás lett kötve. 
Physical integrator struktúra:
-	$1\F$ kondenzátor feszültsége a kondenzátor áramának integrálja osztva $1\F$ kapacitással; azaz numerikusan voltban annyi, mint amennyi Coulomb töltés van a kondenzátoron.
-	$100 \MOhm$ söntellenállás: a Spice szimulátorokban minden csomópontnak szükséges egy egyenáramú úton (ún. DC path) keresztül a földpotenciálra kötődnie, a numerikus számítások elvégzése miatt (a munkaponti számításokban van jelentősége).\\
A memrisztor áramát különböző gerjesztésekre a következő megfigyeléseket tehetjük:
\begin{Anumerate}
	\elem Egyenfeszültségű gerjesztés esetén a memrisztor természetesen Ohmos ellenállásként viselkedik.
	\elem Szinuszos gerjesztésre megfigyelhető a periódikus válaszjel. Ha idő szerinti ábrázolás helyett a feszültség függvényében ábrázoljuk az áramot, a karakterisztikus hiszterézisgörbét kapjuk. A görbe azonban többféleképpen manipulálható. Töltéshordozók mobilitása, feszültségmaplitudó, frekvencia!
	\elem Háromszögjel esetén a szinuszos gerjesztéshez képest a hiszterézis sokkal élesebben vált, ugyanis ekkor a gerjesztő feszültség deriváltja szakaszonként állandó.
\end{Anumerate}

\textit{A memrisztor viselkedésének befolyásolása}

Szép hiszterézis görbét kapunk $0.7 \V$ amplitúdójú, $0.5 \Hz$ frekvenciájú szinuszhullám esetén, a megadott eszközparaméterekkel.
\begin{figure}[H]
	\centering
	\includegraphics[scale=0.3]{memr_07V_05hz}
	\caption{Feszültség-áram időfüggvény 0.7 V amplitudó és 0.5 Hz frekvenciánál)\ref{appendix:kod}}
\end{figure}

\begin{figure}[H]
	\centering
	\includegraphics[scale=0.3]{memr_07V_05hz_hysteresis}
	\caption{Feszültség-áram hiszterézis 0.7 V amplitudó és 0.5 Hz frekvenciánál)\ref{appendix:kod}}
\end{figure}
Ha az amplitúdót nagymértékben növeljük, az adalékolt régió töltéshordozói a nagy elektromos tér miatt gyorsan mozognak. Ekkor hamarabb éri el az eszköz a telítődést, azaz a teljes hossz mentén alacsony ellenállást jelentő állapotot. Ez a hiszterézisben egy éles, függőleges váltásban jelenik meg. Kis amplitudó ($0.3 V$) esetén pedig a gyenge tér alig változtat a memrisztancián, így a görbe alig tér el a lineáristól.
\begin{figure}[H]
	\centering
	\includegraphics[scale=0.3]{memr_1V5_05hz_hyst}
	\caption{Feszültség-áram hiszterézis 1.5 V amplitudó és 0.5 Hz frekvenciánál)\ref{appendix:kod}}
\end{figure}
\begin{figure}[H]
	\centering
	\includegraphics[scale=0.3]{memr_03V_05hz_hysteresis}
	\caption{Feszültség-áram hiszterézis 0.3 V amplitudó és 0.5 Hz frekvenciánál)\ref{appendix:kod}}
\end{figure}
A frekvencia növelésével kevesebb idő jut a rendszer változására, az ionok mozgása így időben korlátozott; szintén a lineárishoz közeli karakterisztika az eredmény. Alacsony frekvencia esetén pedig van idő elérni a telítődési állapotot.
Az ionok mozgékonysága hasonlóan befolyásolja a hiszterézist. Mozgékony ionok gyorsabb telítődést eredményeznek.
\begin{figure}[H]
	\centering
	\includegraphics[scale=0.3]{memr_07V_1hz_hysteresis}
	\caption{Feszültség-áram hiszterézis 0.7 V amplitudó és 1 Hz frekvenciánál)\ref{appendix:kod}}
\end{figure}
Látható tehát, hogy mind a gerjesztés befolyásolásával (az amplitudó és frekvencia modulációjával), mind az eszköz adalékolásával illetve anyagával (mozgékonyabb ionok alkalmazásával) számíthatóan változtathatjuk a memrisztor viselkedését.\smallskip

\textit{A modellezés korlátai és az ablakfüggvény értelmezése}

Mivel ideális eszközt modellezünk matematikai függvények, eljárások alkalmazásával, bizonyos peremfeltételek esetén korlátokba ütközünk, nem reális vagy eltorzult értékeket kapva. Ami fontos, hogy ismerjük a módszerek korlátait, és hogy a számunkra fontos értéktartományokban (példaul a gerjesztés vagy az eszköz paraméterei) a helyes viselkedést adja vissza a modell.\smallskip

Az integrator wind-up jelensége: Az LTSpice modellben, ha nem alkalmazunk ablakfüggvényt (lsd. Joglekar), amint az ionok átérnek a teljes eszközön (D távolság), a fizikai integrátor tovább számolja a töltésmennyiséget a mérő kondenzátoron. Amint megfordul a gerjesztés polaritása, ez a felgyülemlett töltés elkezd kisülni. Túl alacsony frekvencia esetén ez a felgyülemlett irreális nagyságú lesz. Ez az egyszerű Spice modellben olyan, mintha az ionok meghaladnák a memrisztor kiterjedését, így amikor a gerjesztő feszültség polaritása megfordul, nem az eszköz végétől kellene visszaszorulni az ionoknak (ezáltal a nagy ellenállású régió területe rögtön növekedni kezdene), hanem egy nagyobb távolságról, így valótlan értékeket kapunk a váltás után.
Ha a telítődés jelenségét jól akarjuk modellezni, például ablakfüggvényt érdemes használni. Ennek lényege, hogy a telítődés tartományához közelítve egyre kisebb változást ér el a gerjesztés, és az alacsony régió kiterjedése nem haladja meg az eszköz hosszát (a rendszerváltozó – x-szel jelölve a fenti differenciálegyenletekben – deriváltja fokozatosan a nullához tart). Ehhez egy összetettebb .sub fájlt kell alkalmazni, ami ezt figyelembe veszi.
